\documentclass[aspectratio=169,11pt]{beamer}

% =========================
% English Beamer Lecture
% =========================
% pdflatex-compatible setup
\usepackage[utf8]{inputenc}
\usepackage[T1]{fontenc}
\usepackage[english]{babel}
\usepackage{lmodern}
\usepackage{cmap}
\usepackage{amsmath,amssymb}
\usepackage{booktabs}
\usepackage{array}
\usepackage{tabularx}
\usepackage{listings}
\usepackage{xcolor}

\usetheme{Madrid}
\usecolortheme{seahorse}
\setbeamertemplate{navigation symbols}{}
\setbeamertemplate{footline}[frame number]

\definecolor{codebg}{RGB}{248,248,248}
\definecolor{keywordblue}{RGB}{0,70,140}
\definecolor{commentgray}{RGB}{90,90,90}

\lstdefinestyle{sqlstyle}{
  language=SQL,
  basicstyle=\ttfamily\scriptsize,
  keywordstyle=\color{keywordblue}\bfseries,
  commentstyle=\color{commentgray},
  backgroundcolor=\color{codebg},
  frame=single,
  framerule=0.2pt,
  breaklines=true,
  showstringspaces=false,
  columns=fullflexible
}

\title[Data Independence in DBMS]{Lecture: Data Independence in DBMS}
\subtitle{Physical and logical independence in database systems}
\author{Prepared for a DBMS lecture}
\date{\today}

\AtBeginSection[]{
  \begin{frame}{Section Outline}
    \tableofcontents[currentsection]
  \end{frame}
}

\begin{document}

%------------------------------------------------
\begin{frame}
  \titlepage
\end{frame}

%------------------------------------------------
\begin{frame}{Table of Contents}
\begin{columns}[T,totalwidth=\textwidth]
  \column{0.48\textwidth}
  \tableofcontents[sections={1-4},hideallsubsections]
  \column{0.48\textwidth}
  \tableofcontents[sections={5-8},hideallsubsections]
\end{columns}
\end{frame}

%------------------------------------------------
\begin{frame}{Learning Objectives}
After completing this lecture, students should be able to:
\begin{enumerate}
  \item Explain the concept of \textbf{data independence} in DBMS.
  \item Describe the three-level DBMS architecture: \textbf{Internal Level}, \textbf{Conceptual Level}, and \textbf{View Level}.
  \item Distinguish between \textbf{Physical Data Independence} and \textbf{Logical Data Independence}.
  \item Analyze practical examples of each type of data independence.
  \item Explain why data independence reduces maintenance effort and supports long-term system evolution.
\end{enumerate}
\end{frame}

%================================================
\section{Concept of Data Independence}

%------------------------------------------------
\begin{frame}{What Is Data Independence?}
\begin{block}{Definition}
\textbf{Data independence} is the ability to change the schema at one level of a database architecture without affecting the schema at the next higher level.
\end{block}

\vspace{0.3cm}
In simpler terms:
\begin{itemize}
  \item The internal storage or organization of data may change.
  \item Users and application programs can still access data as before.
  \item Application code does not need major changes after internal database optimizations.
\end{itemize}
\end{frame}

%------------------------------------------------
\begin{frame}{Introductory Example}
Suppose a sales system frequently searches orders by customer ID.

\vspace{0.3cm}
The database administrator creates an index to speed up the query:
\begin{itemize}
  \item The internal access path changes.
  \item Users still use the same application interface.
  \item The application still uses the same SQL query.
\end{itemize}

\begin{alertblock}{Key Idea}
The DBMS hides internal database changes from users and application programs.
\end{alertblock}
\end{frame}

%------------------------------------------------
\begin{frame}{Why Is Data Independence Important?}
\begin{columns}[T]
\column{0.33\textwidth}
\begin{block}{Lower Maintenance}
Developers do not need to update applications after every small database change.
\end{block}

\column{0.33\textwidth}
\begin{block}{Greater Flexibility}
The database can be reorganized or optimized internally with limited impact on users.
\end{block}

\column{0.33\textwidth}
\begin{block}{Long-Term Evolution}
The database can be extended as business requirements change without breaking existing systems.
\end{block}
\end{columns}
\end{frame}

%------------------------------------------------
\begin{frame}{Quick Quiz: Concept of Data Independence}
\textbf{Question 1.} What is data independence in DBMS?
\begin{enumerate}[A.]
  \item The ability to change a schema at one level without affecting the next higher level
  \item The ability to delete all data from the system
  \item The ability to change the user interface color
  \item The ability to replace SQL with HTML
\end{enumerate}
\end{frame}

%------------------------------------------------
\begin{frame}{Quick Quiz: Role of Data Independence}
\textbf{Question 2.} Why does data independence reduce maintenance effort?
\begin{enumerate}[A.]
  \item Developers do not need to modify applications after every small database change
  \item Users must manage disk storage themselves
  \item The database can never be changed
  \item A DBMS no longer needs a schema
\end{enumerate}

\vspace{0.2cm}
\textbf{Question 3.} How does data independence support long-term development?
\begin{enumerate}[A.]
  \item It allows the database to be updated and extended with less risk of breaking existing systems
  \item It requires the whole application to be rewritten every year
  \item It prevents all changes to data structure
  \item It works only with text files
\end{enumerate}
\end{frame}

%================================================
\section{Three-Level DBMS Architecture}

%------------------------------------------------
\begin{frame}{Three-Level DBMS Architecture}
To understand data independence, we need to understand the three levels of DBMS architecture:

\vspace{0.4cm}
\begin{center}
\Large
\textbf{View Level}\\[0.2cm]
$\uparrow$\\[0.2cm]
\textbf{Conceptual Level}\\[0.2cm]
$\uparrow$\\[0.2cm]
\textbf{Internal Level}
\end{center}

\vspace{0.2cm}
\begin{alertblock}{Main Idea}
Each level hides details of the lower level to reduce dependency between applications and physical data storage.
\end{alertblock}
\end{frame}

%------------------------------------------------
\begin{frame}{Internal Level}
\begin{block}{Definition}
The \textbf{Internal Level} describes how data is physically stored in the system.
\end{block}

It includes:
\begin{itemize}
  \item Storage files.
  \item Indexes.
  \item Disk space allocation.
  \item Data compression.
  \item Data access methods.
  \item Internal storage structures.
\end{itemize}
\end{frame}

%------------------------------------------------
\begin{frame}{Conceptual Level}
\begin{block}{Definition}
The \textbf{Conceptual Level} describes the logical structure of the database.
\end{block}

It includes:
\begin{itemize}
  \item Tables and fields.
  \item Attributes.
  \item Relationships among tables.
  \item Primary keys and foreign keys.
  \item Integrity constraints.
\end{itemize}
\end{frame}

%------------------------------------------------
\begin{frame}{View Level}
\begin{block}{Definition}
The \textbf{View Level} describes how users and applications see the data.
\end{block}

Examples:
\begin{itemize}
  \item Students can view only their own grades.
  \item Lecturers can view students in their assigned classes.
  \item The academic office can view academic records.
  \item The finance office can view tuition-related data.
\end{itemize}
\end{frame}

%------------------------------------------------
\begin{frame}{Three Levels and Data Independence}
\begin{center}
\begin{tabular}{lll}
\toprule
\textbf{Level} & \textbf{Main Concern} & \textbf{Typical Users} \\
\midrule
View & User-specific data views & End users, applications \\
Conceptual & Logical database design & Designers, developers \\
Internal & Physical storage details & DBA, DBMS engine \\
\bottomrule
\end{tabular}
\end{center}

\vspace{0.3cm}
\begin{alertblock}{Important}
Data independence is achieved by separating these levels and maintaining mappings among them.
\end{alertblock}
\end{frame}

%================================================
\section{Two Types of Data Independence}

%------------------------------------------------
\begin{frame}{Two Types of Data Independence}
\begin{center}
\small
\begin{tabularx}{\textwidth}{>{\bfseries}lXX}
\toprule
Type & Changes Occur At & Should Not Affect \\
\midrule
Physical Data Independence & Internal Level & Conceptual Level and application programs \\
Logical Data Independence & Conceptual Level & View Level and application programs \\
\bottomrule
\end{tabularx}
\end{center}

\vspace{0.3cm}
\begin{block}{Observation}
Physical data independence is usually easier to achieve than logical data independence.
\end{block}
\end{frame}

%================================================
\section{Physical Data Independence}

%------------------------------------------------
\begin{frame}{Physical Data Independence: Definition}
\begin{block}{Definition}
\textbf{Physical Data Independence} is the ability to change how data is physically stored without affecting the logical schema or user applications.
\end{block}

\vspace{0.3cm}
It is related to changes at the \textbf{Internal Level}.

\vspace{0.3cm}
\begin{alertblock}{Main Focus}
Physical data independence focuses on storage, performance, indexing, and access paths.
\end{alertblock}
\end{frame}

%------------------------------------------------
\begin{frame}{Examples of Physical Changes}
The following changes are examples of physical data independence:
\begin{itemize}
  \item Moving database files from one disk to another.
  \item Migrating from HDD to SSD.
  \item Creating indexes to speed up queries.
  \item Compressing data files to save storage space.
  \item Changing access paths or data access methods.
  \item Changing the internal file organization.
\end{itemize}
\end{frame}

%------------------------------------------------
\begin{frame}[fragile]{SQL Example: Creating an Index}
Assume the system frequently searches orders by \texttt{customer\_id}. The DBA creates an index:

\begin{lstlisting}[style=sqlstyle]
CREATE INDEX idx_orders_customer
ON Orders(customer_id);
\end{lstlisting}

The application can still use the same query:

\begin{lstlisting}[style=sqlstyle]
SELECT *
FROM Orders
WHERE customer_id = 101;
\end{lstlisting}

The logical schema of \texttt{Orders} does not change.
\end{frame}

%------------------------------------------------
\begin{frame}{Benefits of Physical Data Independence}
\begin{itemize}
  \item Allows backend optimization without changing user applications.
  \item Improves query performance through indexes and better access paths.
  \item Supports storage and hardware upgrades.
  \item Reduces the need to modify application programs.
  \item Improves long-term maintainability.
\end{itemize}
\end{frame}

%------------------------------------------------
\begin{frame}{Quick Quiz: Physical Data Independence}
\textbf{Question 1.} What is Physical Data Independence?
\begin{enumerate}[A.]
  \item The ability to change physical storage without affecting the logical schema or application
  \item The ability to change the user interface
  \item The ability to delete the whole database
  \item The ability to rename a course
\end{enumerate}

\vspace{0.2cm}
\textbf{Question 2.} Which example illustrates Physical Data Independence?
\begin{enumerate}[A.]
  \item Creating an index to improve query speed without changing the application
  \item Adding an email field to a user form
  \item Changing the website color
  \item Removing a button from the interface
\end{enumerate}
\end{frame}

%================================================
\section{Logical Data Independence}

%------------------------------------------------
\begin{frame}{Logical Data Independence: Definition}
\begin{block}{Definition}
\textbf{Logical Data Independence} is the ability to change the logical schema, that is, the structure and relationships of data, without affecting views or application programs.
\end{block}

\vspace{0.3cm}
It is related to changes at the \textbf{Conceptual Level}.

\vspace{0.3cm}
\begin{alertblock}{Main Focus}
Logical data independence focuses on tables, attributes, relationships, and views.
\end{alertblock}
\end{frame}

%------------------------------------------------
\begin{frame}{Examples of Logical Changes}
The following changes are examples of logical data independence:
\begin{itemize}
  \item Adding a new column to a table.
  \item Removing an unused column.
  \item Creating a new relationship between tables.
  \item Splitting one table into several tables.
  \item Merging multiple tables.
  \item Creating a view to simplify data access.
\end{itemize}
\end{frame}

%------------------------------------------------
\begin{frame}[fragile]{SQL Example: Adding a Column}
Suppose the \texttt{Employees} table initially contains:

\begin{lstlisting}[style=sqlstyle]
employee_id, full_name, department
\end{lstlisting}

The company wants to store employee email addresses:

\begin{lstlisting}[style=sqlstyle]
ALTER TABLE Employees
ADD email VARCHAR(100);
\end{lstlisting}

If old applications use only \texttt{employee\_id}, \texttt{full\_name}, and \texttt{department}, they can continue to work.
\end{frame}

%------------------------------------------------
\begin{frame}[fragile]{Role of Views}
When the logical schema changes, a view can help preserve a stable data interface for applications.

\begin{lstlisting}[style=sqlstyle]
CREATE VIEW EmployeeBasicView AS
SELECT employee_id, full_name, department
FROM Employees;
\end{lstlisting}

Old applications can query \texttt{EmployeeBasicView} without knowing that \texttt{Employees} now contains an extra column.
\end{frame}

%------------------------------------------------
\begin{frame}{Benefits of Logical Data Independence}
\begin{itemize}
  \item Makes application code easier to maintain.
  \item Supports system updates as business requirements evolve.
  \item Reduces the need to rewrite old queries.
  \item Allows the database schema to evolve more flexibly.
  \item Supports multiple views for different user groups.
\end{itemize}
\end{frame}

%------------------------------------------------
\begin{frame}{Quick Quiz: Logical Data Independence}
\textbf{Question 1.} What is Logical Data Independence?
\begin{enumerate}[A.]
  \item The ability to change the logical schema without affecting views or application programs
  \item The ability to change the physical file location
  \item The ability to move data from HDD to SSD
  \item The ability to change the operating system
\end{enumerate}

\vspace{0.2cm}
\textbf{Question 2.} Which example illustrates Logical Data Independence?
\begin{enumerate}[A.]
  \item Adding an \texttt{email} column to \texttt{Employees} while old applications still work
  \item Creating an index on disk
  \item Moving data to an SSD
  \item Compressing a physical data file
\end{enumerate}
\end{frame}

%================================================
\section{Comparison and Integrated Examples}

%------------------------------------------------
\begin{frame}{Physical vs. Logical Data Independence}
\scriptsize
\begin{tabularx}{\textwidth}{lXX}
\toprule
\textbf{Criterion} & \textbf{Physical Data Independence} & \textbf{Logical Data Independence} \\
\midrule
Main focus & Physical data storage & Logical data structure \\
Related level & Internal Schema & Conceptual Schema \\
Typical changes & Files, indexes, disks, compression, access paths & Tables, columns, relationships, views \\
Impact on applications & Should not affect application programs & Harder to avoid application impact \\
Main goal & Improve performance and storage & Evolve database design \\
Difficulty & Usually easier & Usually harder \\
Example & Moving files or adding an index & Adding or removing a column \\
\bottomrule
\end{tabularx}
\end{frame}

%------------------------------------------------
\begin{frame}[fragile]{Integrated Example: Sales System}
Consider a sales management system.

\vspace{0.2cm}
\textbf{Physical Data Independence:} The DBA creates an index on \texttt{customer\_id} in \texttt{Orders}.
\begin{lstlisting}[style=sqlstyle]
CREATE INDEX idx_customer_id
ON Orders(customer_id);
\end{lstlisting}

The old query still works:
\begin{lstlisting}[style=sqlstyle]
SELECT *
FROM Orders
WHERE customer_id = 1001;
\end{lstlisting}
\end{frame}

%------------------------------------------------
\begin{frame}[fragile]{Integrated Example: Logical Change}
\textbf{Logical Data Independence:} The company wants to store customer email addresses.

\begin{lstlisting}[style=sqlstyle]
ALTER TABLE Customers
ADD email VARCHAR(100);
\end{lstlisting}

Old reports that use only \texttt{customer\_id} and \texttt{customer\_name} can still work.

\vspace{0.2cm}
A view can preserve the old structure:
\begin{lstlisting}[style=sqlstyle]
CREATE VIEW CustomerBasicView AS
SELECT customer_id, customer_name
FROM Customers;
\end{lstlisting}
\end{frame}

%================================================
\section{Review Questions and Exercises}

%------------------------------------------------
\begin{frame}{Review Questions: Multiple Choice I}
\textbf{Question 1.} What is data independence?
\begin{enumerate}[A.]
  \item The ability to change a schema at one level without affecting the next higher level
  \item The ability to delete data without control
  \item The ability to change the interface color
  \item The ability to replace a database with an image file
\end{enumerate}

\vspace{0.2cm}
\textbf{Question 2.} How many main types of data independence are there?
\begin{enumerate}[A.]
  \item 1
  \item 2
  \item 3
  \item 4
\end{enumerate}
\end{frame}

%------------------------------------------------
\begin{frame}{Review Questions: Multiple Choice II}
\textbf{Question 3.} Physical Data Independence is mainly related to which level?
\begin{enumerate}[A.]
  \item Internal Level
  \item View Level
  \item User Interface
  \item Application Screen
\end{enumerate}

\vspace{0.2cm}
\textbf{Question 4.} Logical Data Independence is mainly related to which level?
\begin{enumerate}[A.]
  \item Internal Level
  \item Conceptual Level
  \item Disk Level
  \item Hardware Level
\end{enumerate}
\end{frame}

%------------------------------------------------
\begin{frame}{Review Questions: Multiple Choice III}
\textbf{Question 5.} Which example illustrates Physical Data Independence?
\begin{enumerate}[A.]
  \item Creating a new index without changing the application
  \item Adding an email field to the user interface
  \item Changing the website color
  \item Removing a button from the screen
\end{enumerate}

\vspace{0.2cm}
\textbf{Question 6.} Which example illustrates Logical Data Independence?
\begin{enumerate}[A.]
  \item Adding a new column to a table while old applications still work
  \item Moving data to an SSD
  \item Creating a physical file
  \item Changing the disk block size
\end{enumerate}
\end{frame}

%------------------------------------------------
\begin{frame}{Short-Answer Questions}
\begin{enumerate}
  \item Explain the concept of data independence in DBMS.
  \item Distinguish between Physical Data Independence and Logical Data Independence.
  \item Why can Physical Data Independence improve performance without affecting applications?
  \item Why is Logical Data Independence usually harder to achieve than Physical Data Independence?
  \item Explain the role of views in supporting Logical Data Independence.
\end{enumerate}
\end{frame}

%------------------------------------------------
\begin{frame}{Applied Exercises}
\begin{enumerate}
  \item A DBA moves data from HDD to SSD and creates an index to speed up queries. The application does not change. Identify the type of data independence and explain.
  \item An HR system adds an \texttt{email} column to \texttt{Employees}. Old reports still show only employee ID, name, and department. Identify the type of data independence and explain.
  \item A sales system changes the structure of \texttt{Customers}, but old applications still need \texttt{customer\_id} and \texttt{customer\_name}. Propose a view-based solution.
  \item Give one student management system example for each type of data independence.
\end{enumerate}
\end{frame}

%================================================
\section{Summary}

%------------------------------------------------
\begin{frame}{Lesson Summary}
\begin{itemize}
  \item Data independence allows schema changes at one level without affecting the next higher level.
  \item DBMS architecture usually has three levels: Internal, Conceptual, and View.
  \item Physical Data Independence allows changes in physical storage without affecting the logical schema or applications.
  \item Logical Data Independence allows changes in the logical schema without affecting views or applications.
  \item Physical Data Independence is usually easier to achieve than Logical Data Independence.
  \item Views can help hide logical schema changes and keep applications stable.
\end{itemize}
\end{frame}

%------------------------------------------------
\begin{frame}{Key Terms}
\begin{columns}[T]
\column{0.48\textwidth}
\begin{itemize}
  \item Data Independence
  \item Physical Data Independence
  \item Logical Data Independence
  \item DBMS
  \item Internal Level
  \item Conceptual Level
  \item View Level
\end{itemize}

\column{0.48\textwidth}
\begin{itemize}
  \item Internal Schema
  \item Conceptual Schema
  \item View
  \item Index
  \item Schema
  \item Database Administrator
  \item Application Program
\end{itemize}
\end{columns}
\end{frame}

%------------------------------------------------
\begin{frame}{Closing Thought}
\begin{block}{Core Message}
Data independence is one of the key reasons why DBMSs are more powerful and maintainable than direct file-based data management.
\end{block}

\vspace{0.4cm}
A well-designed database system separates:
\begin{itemize}
  \item how data is physically stored,
  \item how data is logically organized,
  \item and how users or applications see the data.
\end{itemize}
\end{frame}

\end{document}
